\documentclass[10pt,a4paper]{amsart}
\usepackage{amsmath,amssymb,amsthm,graphicx}
\usepackage{verbatim}
\newenvironment{reference}{\comment}{\endcomment}
%\newenvironment{reference}{}{}
\newenvironment{slogan}{\comment}{\endcomment}
\newenvironment{history}{\comment}{\endcomment}

% For commutative diagrams we use Xy-pic
\usepackage[all]{xy}

% We use 2cell for 2-commutative diagrams.
\xyoption{2cell}
% Theorem environments.
%
\theoremstyle{plain}
\newtheorem{theorem}[subsection]{Theorem}
\newtheorem{proposition}[subsection]{Proposition}
\newtheorem{lemma}[subsection]{Lemma}

\theoremstyle{definition}
\newtheorem{definition}[subsection]{Definition}
\newtheorem{example}[subsection]{Example}
\newtheorem{exercise}[subsection]{Exercise}
\newtheorem{situation}[subsection]{Situation}

\theoremstyle{remark}
\newtheorem{remark}[subsection]{Remark}
\newtheorem{remarks}[subsection]{Remarks}

\numberwithin{equation}{subsection}

% Macros
%
\def\lim{\mathop{\mathrm{lim}}\nolimits}
\def\colim{\mathop{\mathrm{colim}}\nolimits}
\def\Spec{\mathop{\mathrm{Spec}}}
\def\Hom{\mathop{\mathrm{Hom}}\nolimits}
\def\Ext{\mathop{\mathrm{Ext}}\nolimits}
\def\SheafHom{\mathop{\mathcal{H}\!\mathit{om}}\nolimits}
\def\SheafExt{\mathop{\mathcal{E}\!\mathit{xt}}\nolimits}
\def\Sch{\mathit{Sch}}
\def\Mor{\mathop{\mathrm{Mor}}\nolimits}
\def\Ob{\mathop{\mathrm{Ob}}\nolimits}
\def\Sh{\mathop{\mathit{Sh}}\nolimits}
\def\NL{\mathop{N\!L}\nolimits}
\def\CH{\mathop{\mathrm{CH}}\nolimits}
\def\proetale{{pro\text{-}\acute{e}tale}}
\def\etale{{\acute{e}tale}}
\def\QCoh{\mathit{QCoh}}
\def\Ker{\mathop{\mathrm{Ker}}}
\def\Im{\mathop{\mathrm{Im}}}
\def\Coker{\mathop{\mathrm{Coker}}}
\def\Coim{\mathop{\mathrm{Coim}}}

% Boxtimes
%
\DeclareMathSymbol{\boxtimes}{\mathbin}{AMSa}{"02}

%
% Macros for moduli stacks/spaces
%
\def\QCohstack{\mathcal{QC}\!\mathit{oh}}
\def\Cohstack{\mathcal{C}\!\mathit{oh}}
\def\Spacesstack{\mathcal{S}\!\mathit{paces}}
\def\Quotfunctor{\mathrm{Quot}}
\def\Hilbfunctor{\mathrm{Hilb}}
\def\Curvesstack{\mathcal{C}\!\mathit{urves}}
\def\Polarizedstack{\mathcal{P}\!\mathit{olarized}}
\def\Complexesstack{\mathcal{C}\!\mathit{omplexes}}
% \Pic is the operator that assigns to X its picard group, usage \Pic(X)
% \Picardstack_{X/B} denotes the Picard stack of X over B
% \Picardfunctor_{X/B} denotes the Picard functor of X over B
\def\Pic{\mathop{\mathrm{Pic}}\nolimits}
\def\Picardstack{\mathcal{P}\!\mathit{ic}}
\def\Picardfunctor{\mathrm{Pic}}
\def\Deformationcategory{\mathcal{D}\!\mathit{ef}}
\usepackage{cleveref}
\crefname{originaldefinition}{original definition}{original definitions}
\Crefname{originaldefinition}{Original definition}{Original definitions}
\makeatletter
\newlabel{definition-filtered}{{0121}{source}{}{sourcepointer-1}{}}
\newlabel{definition-filtered@cref}{{[originaldefinition][0][]0121}{[1][0][]source}{}{}{}}
\makeatother
\crefname{originallemma}{original lemma}{original lemmas}
\Crefname{originallemma}{Original lemma}{Original lemmas}
\makeatletter
\newlabel{lemma-biregular-ss-converges}{{012W}{source}{}{sourcepointer-2}{}}
\newlabel{lemma-biregular-ss-converges@cref}{{[originallemma][0][]012W}{[1][0][]source}{}{}{}}
\makeatother
\title{The class of a finitely supported spectral sequence in K-zero equals the class of its cohomology}
\author{The Stacks Project Authors}
\date{stacks-project:homology.tex:local-snapshot-20261012; GNU FDL 1.2 or later}
\begin{document}
\maketitle
\noindent\textit{Editorial note.} All delivered mathematical text is the author's own, copied from the pinned Stacks Project chapter file homology.tex; no statement or proof step was rewritten, reordered or removed.\par\medskip
\noindent\textit{Editorial note.} The proof is the author's complete argument; the closing sentence 'and the result is clear' follows a displayed finite filtration whose graded pieces were identified in the cited convergence lemma, so it closes the argument rather than deferring a step.\par\medskip
\noindent\textit{Editorial note.} Every labelled reference occurring inside the counted spans is carried in assembly.external\_references as a non-load-bearing navigation pointer (same-chapter ones at their definition site, cross-chapter ones at the invocation inside the counted span), so no load-bearing mathematics is replaced by a stub. No bibliography citation occurs inside the counted spans.\par\medskip
\noindent\textit{Editorial note.} The same-name check for this card scanned every staged Stacks chapter of this repository -- all sources/stacks-*/pinned-original-source.tex files, the card's own chapter excluded -- and the labels found are recorded in dedup.same\_name\_cross\_chapter on the card.\par\medskip
\noindent\textit{Editorial note.} The preamble is the chapter's own native preamble excerpt already used for this source: the comment-type reference, slogan and history environments (native preamble lines 8-13), the Xy-pic setup (native preamble lines 14-18) and the native theorem declarations with the mathematical macros (native preamble lines 154-217, the actual last line of the file).\par\medskip
\noindent\textit{Allowed standard prerequisites (editorial).}\begin{itemize}
\item Filtered complexes and the spectral sequence attached to one. Same chapter, Definition definition-filtered together with the construction of the spectral sequence of a filtered complex. Supplies the filtered complex, the finiteness of the filtration and the pages of the spectral sequence on which the argument operates.
\item Convergence of the spectral sequence of a finitely filtered complex. Same chapter, Lemma lemma-biregular-ss-converges. Supplies that each cohomology object has a finite filtration whose graded pieces are the limiting terms, which is what turns the alternating sums into the displayed identity.
\item The zeroth K-group of an abelian category. Same chapter, Definition definition-K-zero. Supplies the group in which the displayed identity is taken and in which alternating sums of classes are computed.
\end{itemize}\medskip
\ifdefined\hypertarget\hypertarget{sourcepointer-1}{}\fi
\noindent\textit{Source reference.} Original definition 0121: sources/stacks-homology/pinned-original-source.tex, source lines 4884--4884. Navigation pointer invoked at proof line 6582; the counted span contains the pointer only and the cited result is not reproduced inside it, so the pointer is registered as a non-load-bearing navigation stub in the established form for Stacks material. No bibliography citation is involved.\par\medskip
\ifdefined\hypertarget\hypertarget{sourcepointer-2}{}\fi
\noindent\textit{Source reference.} Original lemma 012W: sources/stacks-homology/pinned-original-source.tex, source lines 6579--6579. Navigation pointer invoked at proof line 6666; the counted span contains the pointer only and the cited result is not reproduced inside it, so the pointer is registered as a non-load-bearing navigation stub in the established form for Stacks material. No bibliography citation is involved.\par\medskip
\begin{lemma}
\label{lemma-biregular-ss-relation-in-K0}
Let $\mathcal{A}$ be an abelian category. Let $(K^\bullet, F)$ be a
filtered complex of $\mathcal{A}$. Assume that the filtration on each $K^n$
is finite (see Definition \ref{definition-filtered}) and that for some
$r$ we have only a finite number of nonzero $E_r^{p, q}$. Then
only a finite number of $H^n(K^\bullet)$ are nonzero and we have
$$
\sum (-1)^n[H^n(K^\bullet)] = \sum (-1)^{p + q} [E_r^{p, q}]
$$
in $K_0(\mathcal{A}')$ where $\mathcal{A}'$ is the smallest weak
Serre subcategory of $\mathcal{A}$ containing the objects
$E_r^{p, q}$.
\end{lemma}

\begin{proof}
Denote $E_r^{even}$ and $E_r^{odd}$ the even and odd part of $E_r$
defined as the direct sum of the $(p, q)$ components with $p + q$ even
and odd. The differential $d_r$ defines maps
$\varphi : E_r^{even} \to E_r^{odd}$ and $\psi : E_r^{odd} \to E_r^{even}$
whose compositions either way give zero.
Then we see that
\begin{align*}
[E_r^{even}] - [E_r^{odd}] & =
[\Ker(\varphi)] + [\Im(\varphi)] - [\Ker(\psi)] - [\Im(\psi)] \\
& =
[\Ker(\varphi)/\Im(\psi)] - [\Ker(\psi)/\Im(\varphi)] \\
& =
[E_{r + 1}^{even}] - [E_{r + 1}^{odd}]
\end{align*}
Note that all the intervening objects are in the smallest Serre
subcategory containing the objects $E_r^{p, q}$.
Continuing in this manner we see that we can increase $r$ at will.
Since there are only a finite number of pairs $(p, q)$ for which
$E_r^{p, q}$ is nonzero, a property which is inherited by
$E_{r + 1}, E_{r + 2}, \ldots$, we see that we may assume that $d_r = 0$.
At this stage we see that $H^n(K^\bullet)$ has a finite filtration
(Lemma \ref{lemma-biregular-ss-converges}) whose graded pieces
are exactly the $E_r^{p, n - p}$ and the result is clear.
\end{proof}

\end{document}
